\section{Introduction}

Cortical neurons receive thousands of synaptic inputs distributed across elaborate dendritic trees. A growing body of evidence suggests that dendrites are not passive conduits for synaptic signals, but active computational elements that perform local nonlinear operations on spatially clustered inputs \citep{poirazi2003,branco2010branch,major2013}. Theoretical and experimental work has shown that individual dendritic branches can act as semi-independent computational subunits, generating local regenerative events such as NMDA spikes \citep{schiller2000} and calcium plateaus \citep{larkum2022} that depend critically on the spatial arrangement of co-active synapses.

If dendritic branches function as local integrators, then the spatial organization of synaptic inputs along the tree is not merely an anatomical detail but a key determinant of neuronal computation. This raises a central question: are synapses from the same presynaptic neuron spatially clustered on dendrites? Such clustering would create conditions favorable for local dendritic integration of correlated inputs, consistent with Hebbian models of synaptic organization \citep{kleindienst2011,takahashi2012}.

Several lines of evidence point toward input-specific spatial organization. Calcium imaging studies have revealed that functionally co-tuned synapses tend to cluster on the same dendritic branches \citep{wilson2016,jia2010}, and electron microscopy reconstructions have demonstrated structured connectivity between identified neuronal populations \citep{druckmann2014,kasthuri2015}. However, these studies have generally been limited to small numbers of identified inputs or to functional rather than anatomical measures of clustering. The recent availability of large-scale connectomic datasets, particularly the MICrONS cortical connectome \citep{microns2025}, now makes it possible to test for input-specific spatial organization across hundreds of presynaptic partners on individual neurons.

A key challenge in detecting partner-specific clustering is controlling for confounds that can produce apparent spatial structure. Synapse density varies systematically with distance from the soma, and this first-order gradient can create the appearance of clustering if not properly controlled. Previous statistical frameworks for detecting spatial organization on dendrites, such as the Neurostat multiscale analysis \citep{neurostat}, have used uniform null models that do not account for distance-dependent density variations. Furthermore, testing hundreds of partners per neuron requires careful multiple testing correction and validation against empirical null distributions.

Here we address these challenges with a two-stage analytical framework. First, we construct a soma-distance null model that accounts for known distance-dependent density gradients, and show that spatial structure persists beyond this first-order organization in all 12 neurons analyzed. Second, we develop a per-partner clustering test using distance-matched permutation nulls that preserve each partner's soma-distance profile, enabling rigorous detection of partner-specific spatial compactness. We validate our findings with a label-shuffle control that confirms the signal is tied to presynaptic identity, and we characterize the clustering as branch-local, operating at a scale of approximately 86~$\mu$m. Our results provide quantitative evidence that input-specific dendritic clustering is a common feature of cortical neurons, consistent with models of spatially organized synaptic connectivity that support local dendritic computation.
